\section{The three questions and their attribution}\label{sec:introduction}

A set \(A\subseteq\N=\{1,2,\ldots\}\) has \emph{Property P} if there
are no pairwise distinct elements \(a,b,c\in A\) such that
\[
 a<b,\qquad a<c,\qquad a\mid b+c.
\]
In particular, equality of the two larger summands is permitted by the
formulation and does not constitute a forbidden triple. Write
\(A(N)=|A\cap\{1,\ldots,N\}|\).
Erd\H{o}s Problem~12 asks three separate questions about the counting
function and reciprocal sum of such an infinite set.

\begin{theorem}[Part (i), AlphaProof Nexus]\label{thm:part-i}
There is an infinite set \(A\) with Property P such that
\[
 \liminf_{N\to\infty}\frac{A(N)}{\sqrt N}>0.
\]
\end{theorem}

\begin{theorem}[Part (ii), AlphaProof Nexus]\label{thm:part-ii}
There is no absolute \(c>0\) such that every infinite set \(A\) with
Property P satisfies \(A(N)<N^{1-c}\) for infinitely many \(N\).
More precisely, for each \(0<c<1\) there are an infinite set \(A_c\)
with Property P and an integer \(N_c\) such that
\begin{equation}\label{eq:parameterwise-dense}
 A_c(N)\ge N^{1-c}\qquad(N\ge N_c).
\end{equation}
The set and the threshold may depend on \(c\).
\end{theorem}

\begin{theorem}[Part (iii), with a uniform bound]\label{thm:main}
There is an absolute constant \(C\) such that every set
\(A\subseteq\N\) with Property P satisfies
\begin{equation}\label{eq:main}
 \hrec(A):=\sum_{a\in A}\frac1a\le C.
\end{equation}
The same constant applies to finite and infinite sets.
\end{theorem}

The answers to the three numbered questions are therefore, respectively,
\emph{yes}, \emph{no}, and \emph{yes}. The first two answers are prior
results of the Google DeepMind AlphaProof Nexus work of Tsoukalas,
Kovsharov, Shirobokov et al.\ \cite{TKS}. They are included here with
attribution, exposition, and adaptations of the released Lean proofs
\cite{NexusI,NexusII}; they are not contributions first obtained in this
manuscript. The reciprocal-sum argument of Theorem~\ref{thm:main} is the
argument developed in the original manuscript by the four named authors.
The present version brings the three statements into one account and one
checked Lean source file.

The quantifier order in \eqref{eq:parameterwise-dense} matters:
\(\forall c\,\exists A_c\,\exists N_c\). We use precisely this
parameterwise assertion. The construction below does not assert that one
fixed set works for every \(c\). The lower bounds in part (ii) are
compatible with the uniform reciprocal bound in part (iii): the onset
\(N_c\) and the set itself vary with the exponent.

Erd\H{o}s and S\'ark\"ozy \cite{ES} proved that every infinite set with
Property P has density zero, and formulated the reciprocal-sum problem
and its stronger uniform version. For the counting questions,
Section~\ref{sec:counting} follows the congruence-block constructions in
the released AlphaProof Nexus Lean files. Those are the concrete proof
instances incorporated in the companion formalization. The constructions
are presented here directly, rather than by relying on admitted statements
in the Formal Conjectures catalog.

The proof of part (iii) has three main steps. A quantitative
incompatibility between separated blocks of rough elements is combined
with sieve estimates to control prime fibres carrying inherited congruence
restrictions. We then represent reciprocal sums by random prime-factor
prefix weights. An adapted rule for selecting the restrictions gives a
uniform total expected weight. The restriction estimate is a special case
of Bera--Viswanadham \cite{BV}; Section~\ref{sec:restriction} includes a
direct proof in the form needed here. The Mertens prefix weights are
classical; see Alexeev et al.\ \cite{AB}. Earlier work of Plasma AI
\cite{Plasma} contains the reduction to odd primitive sets, restricted
reciprocal-sum estimates, and the problem of controlling accumulated
congruence restrictions. These attributions apply to the reciprocal-sum
argument, independently of the provenance of parts (i) and (ii).

\subsection*{Notation and conventions}
All logarithms are natural unless a base is indicated.
We write \(\e(t)=\exp(2\pi i t)\), and \(\norm{t}_{\T}\) for distance
from a real number to the nearest integer. The constants implicit in
\(\ll\) and \(\asymp\) are absolute unless a dependence is indicated.
An integer is \emph{\(D\)-rough} if no prime factor is at most \(D\),
and \emph{\(D\)-smooth} if every prime factor is at most \(D\).
The integer 1 satisfies both conventions. We write \(\Pmax(g)\) for
the largest prime factor of \(g>1\), with \(\Pmax(1)=1\).
A set is \emph{primitive} if no two distinct members divide one another.
Property P is inherited by subsets and by division by a common divisor.
The paper uses positive integers; the Lean statements use natural numbers,
with \(1/0=0\). The constructions below consist of positive integers.

\section{The counting constructions: parts (i) and (ii)}\label{sec:counting}

This section is an attributed exposition of the released proofs
\cite{NexusI,NexusII}. We isolate the common congruence argument and then
make the two choices of scale. Omitting finitely many initial blocks in
the exposition simplifies the estimates and does not affect the asserted
asymptotic statements. The Lean implementations verify explicit initial
bounds as well.

\begin{lemma}[Congruence separation]\label{lem:congruence-separation}
Let \(B_j\) be finite sets of positive integers, with every element of
\(B_i\) smaller than every element of \(B_j\) when \(i<j\). Suppose that
there are pairwise distinct primes \(q_j>2\) such that elements of
\(B_j\) are \(0\pmod{q_j}\) and \(1\pmod{q_i}\) for every \(i<j\).
If each \(B_j\) lies in an interval \([L_j,U_j]\) with
\(2U_j<3L_j\), then \(\bigcup_j B_j\) has Property P.
\end{lemma}
\begin{proof}
Suppose \(a\in B_i\), \(a<b,c\), and \(a\mid b+c\).
If either larger element lies in a later block, then \(b+c\) is
\(1\) or \(2\pmod{q_i}\), whereas \(a\equiv0\pmod{q_i}\).
This is impossible because \(q_i>2\). Otherwise both larger elements
lie in \(B_i\), and
\(2a<b+c\le2U_i<3L_i\le3a\), which precludes divisibility.
\end{proof}

\subsection{A positive square-root lower limit}

Choose increasing primes \(q_0,q_1,\ldots\), starting with \(q_0=5\),
so that \(q_{j+1}\le2q_j+2\). Bertrand's postulate gives such a choice,
and induction yields \(q_j+2\le7\cdot2^j\). Put
\[
 K_j=729\prod_{i<j}q_i,\qquad P_j=K_jq_j,\qquad M_j=P_j^8,
 \qquad T_j=\left\lfloor\frac{P_{j+1}^4}{2}\right\rfloor.
\]
The Chinese remainder theorem supplies \(r_j\in[0,P_j)\) such that
\(r_j\equiv1\pmod{K_j}\) and \(r_j\equiv0\pmod{q_j}\).
Define
\[
 B_j=\{M_j+r_j+tP_j:0\le t<T_j\}.
\]
Then \(|B_j|=T_j\). Also \(P_{j+1}=P_jq_{j+1}\), while
\(P_j\ge3645\cdot5^j\). Consequently
\[
 \frac{q_{j+1}^4}{P_j^3}\longrightarrow0,\qquad
 \frac{T_jP_j}{M_j}\longrightarrow0.
\]
Thus, after removing finitely many initial blocks, each \(B_j\) lies
in \([M_j,(7/5)M_j]\). The scales are strictly separated because
\(M_{j+1}=M_jq_{j+1}^8\). The residues of \(B_j\) have the form in
Lemma~\ref{lem:congruence-separation}, since \(q_i\mid K_j\) for
\(i<j\). Their union \(A\) therefore has Property P.

For large \(j\),
\[
 T_j=\tfrac12\sqrt{M_{j+1}}+O(1).
\]
If \(N\) falls inside the short interval containing \(B_j\), the
preceding complete block contributes
\(T_{j-1}=\tfrac12\sqrt{M_j}+O(1)\) elements. If \(N\) falls in the
following gap, block \(B_j\) contributes
\(T_j=\tfrac12\sqrt{M_{j+1}}+O(1)\) elements. These two cases give
\[
 A(N)\ge\tfrac13\sqrt N
\]
for every sufficiently large \(N\).
To verify that the lower limit is finite as well, note that \(r_j>0\),
so \(A(M_j)=\sum_{i<j}|B_i|\), up to the omitted finite prefix. Since
\(P_{i+1}\le P_j5^{-(j-i-1)}\),
\[
 A(M_j)\le\frac12 P_j^4\sum_{k\ge0}5^{-4k}<P_j^4=\sqrt{M_j}.
\]
This proves Theorem~\ref{thm:part-i}. In particular, the lower limit is
bounded away from zero and bounded above; this agrees with the real-valued
lower-limit convention used by the Lean statement.

\subsection{Failure of a universal power saving}

Fix \(0<c<1\). Choose distinct odd primes \(p_j\) with
\(p_j\le2^{j+2}\), and put \(Q_j=\prod_{i\le j}p_i\). Then
\(Q_j\le2^{(j+2)^2}\). Choose an integer \(a\ge4\) with \(ac\ge3\),
and define
\[
 M_j=2^{a(j+2a)^2}.
\]
Choose a residue \(s_j\pmod{Q_j}\) that is \(0\pmod{p_j}\) and
\(1\pmod{p_i}\) for \(i<j\), and set
\[
 B_j=\{n\in[10M_j,14M_j]:n\equiv s_j\pmod{Q_j}\}.
\]
The intervals are disjoint and ordered because
\(14M_j<10M_{j+1}\). They satisfy the strict width inequality in
Lemma~\ref{lem:congruence-separation}. Once again their union \(A_c\)
has Property P, and it is infinite because the blocks are nonempty for
all sufficiently large \(j\).

Counting a residue class in an interval gives
\[
 |B_j|\ge\frac{4M_j}{Q_j}-1.
\]
The logarithm to base 2 of \(4M_j/Q_j\) is at least
\(2+a(j+2a)^2-(j+2)^2\), whereas
\[
 \log_2\bigl((14M_{j+1})^{1-c}\bigr)
 =(1-c)\bigl(\log_2 14+a(j+1+2a)^2\bigr).
\]
The difference has leading term \((ac-1)j^2\), whose coefficient is
positive. Therefore, for all sufficiently large \(j\),
\[
 |B_j|\ge(14M_{j+1})^{1-c}.
\]
For \(14M_j\le N<14M_{j+1}\), the complete block \(B_j\) is already
counted in \(A_c(N)\). This proves \eqref{eq:parameterwise-dense}.
If a universal saving \(d>0\) existed, apply the construction with
\(c=\min(d,1/2)\). For sufficiently large \(N\ge1\),
\(A_c(N)\ge N^{1-c}\ge N^{1-d}\), contradicting the required
infinitely many strict inequalities. This proves
Theorem~\ref{thm:part-ii}.

